%%%%%%%%%%%%%%%%%%%%%%%%%%%%%%%%%%%%%%%%%%%%%%%%%%%%%%%%%%%%%%%%%%%%%%%%%%%%%%%
% CHAPTER 68: DEFORMATION QUANTIZATION EXAMPLES
%%%%%%%%%%%%%%%%%%%%%%%%%%%%%%%%%%%%%%%%%%%%%%%%%%%%%%%%%%%%%%%%%%%%%%%%%%%%%%%

\chapter{Deformation Quantization Examples}
\label{chap:deformation-examples}


\section{\texorpdfstring{$\Pinf$ Structures}{Pinf Structures}}
\label{sec:pinf-structures}
\begin{definition}[$\Pinf$-Chiral Algebra]\label{def:pinf-chiral-detailed}
A \textbf{$\Pinf$-chiral algebra} is a chiral algebra $\mathcal{A}$ equipped with:
\begin{enumerate}[label=(\roman*)]
\item An $\Einf$-chiral algebra structure (commutative OPE).
\item A compatible $\Linf$-chiral Lie structure $\{-, -\}$ (Poisson bracket).
\item The Leibniz rule: $\{a, b \cdot c\} = \{a, b\} \cdot c + b \cdot \{a, c\}$.
\end{enumerate}
Equivalently, $\mathcal{A}$ is an algebra over the chiral Poisson operad $\chirPois$.
\end{definition}

\begin{example}[Classical Affine Poisson]\label{ex:classical-affine}
Let $\mathfrak{g}^*$ be the dual of a Lie algebra with Kirillov--Kostant Poisson structure:
\[
\{f, g\}(x) = x([df_x, dg_x])
\]
The loop algebra version $L\mathfrak{g}^* = \mathfrak{g}^* \otimes \bC[t, t^{-1}]$ has a
$\Pinf$-chiral structure:
\[
\{J^a(z), J^b(w)\} = f^{ab}_c J^c(w) \delta(z - w)
\]
This is the classical limit ($k \to \infty$) of the Kac--Moody OPE.
\end{example}


\subsection{Coisson Algebras}

\begin{definition}[Coisson Algebra]\label{def:coisson}
A \textbf{coisson algebra} is a commutative chiral algebra $\mathcal{A}$ with a 
compatible chiral Lie cobracket:
\[
\delta: \mathcal{A} \to \mathcal{A} \chirtensor \mathcal{A}
\]
satisfying the co-Leibniz rule.
\end{definition}

\begin{theorem}[Coisson = $(\chirPois)^c$-Coalgebra; \ClaimStatusProvedElsewhere]\label{thm:coisson-coalgebra}
Coisson algebras are exactly coalgebras over the Koszul dual cooperad $(\chirPois)^c$.
\end{theorem}


\section{Quantization Formulas}
\label{sec:quantization-formulas}

\subsection{Deformation Quantization Setup}

\begin{construction}[Quantization Map]\label{constr:quantization}
A \textbf{deformation quantization} of a $\Pinf$-chiral algebra $(\mathcal{A}_0, \cdot, \{-,-\})$ 
is an $\Eone$-chiral algebra $(\mathcal{A}_\hbar, \star)$ such that:
\begin{enumerate}[label=(\roman*)]
\item $\mathcal{A}_\hbar = \mathcal{A}_0[[\hbar]]$ as a vector space.
\item $a \star b = a \cdot b + \hbar B_1(a, b) + \hbar^2 B_2(a, b) + \cdots$.
\item $a \star b - b \star a = \hbar \{a, b\} + O(\hbar^2)$.
\end{enumerate}
\end{construction}

\begin{theorem}[Formality for $\Pinf$-Chiral; \ClaimStatusProvedElsewhere]\label{thm:pinf-formality}
Every $\Pinf$-chiral algebra admits a deformation quantization to an $\Eone$-chiral
algebra. The quantization is unique up to gauge equivalence.
\end{theorem}

\begin{proof}
The result follows from the formality of the $E_2$-operad (Kontsevich 1999,
Tamarkin 2003) via the factorization algebra formalism.  Specifically:

\textbf{Step 1: Formality.}
The Francis--Gaitsgory formality theorem
(Theorem~\ref{thm:chiral-formality}, \ClaimStatusProvedElsewhere)
establishes an $L_\infty$ quasi-isomorphism between chiral polyvector
fields and chiral Hochschild cochains.  This guarantees that the
deformation problem is unobstructed: the obstruction class in
$H^3_{\chirPois}(\mathcal{A}_0; \mathcal{A}_0)$ vanishes by formality.

\textbf{Step 2: Existence.}
By formality, a Poisson structure $\pi \in H^2_{\chirPois}$ lifts to a
Maurer--Cartan element $\alpha = \hbar\pi + O(\hbar^2)$ in the
Hochschild cochain complex.  The resulting star product is:
\[
a \star b = \sum_{n \geq 0} \hbar^n \sum_{\Gamma \in G_{n,2}} w_\Gamma B_\Gamma(a, b)
\]
where $G_{n,2}$ are admissible graphs, $w_\Gamma$ are Kontsevich weights
(integrals over $\FM_n(\mathbb{C})$), and $B_\Gamma$ are bidifferential
operators.  Associativity follows from Stokes' theorem on the boundary
strata of $\FM_n(\mathbb{C})$.

\textbf{Step 3: Uniqueness.}
Two quantizations correspond to gauge-equivalent Maurer--Cartan elements
in the deformation complex (Theorem~\ref{thm:mc-quantization}).  The
gauge group acts transitively on the fiber over a fixed first-order
deformation by standard obstruction theory: $H^1$ controls the
infinitesimal gauge transformations, and the formality quasi-isomorphism
ensures no higher obstructions to gauge equivalence.
\end{proof}


\subsection{Explicit Star Product}

\begin{computation}[Star Product through Order $\hbar^2$]\label{comp:star-order-2}
For a $\Pinf$-chiral algebra with Poisson bracket $\{a, b\}$:

\textbf{Order $\hbar^0$:} $B_0(a, b) = a \cdot b$.

\textbf{Order $\hbar^1$:} $B_1(a, b) = \frac{1}{2}\{a, b\}$.

\textbf{Order $\hbar^2$:}  In local coordinates with Poisson bivector $\pi^{ij}$:
\[
B_2(a, b) = \frac{1}{8}\pi^{ij}\pi^{kl}\partial_i\partial_k a \cdot \partial_j\partial_l b + \frac{1}{12}\pi^{ij}(\partial_j\pi^{kl})(\partial_i\partial_k a \cdot \partial_l b - \partial_k a \cdot \partial_i\partial_l b)
\]
The first term corresponds to the ``double Poisson'' Kontsevich graph (both internal vertices connecting directly to the two external vertices); the second involves $\partial\pi$ and vanishes for constant Poisson structures.
\end{computation}


\section{Obstructions and Anomalies in Examples}
\label{sec:obstructions}

\subsection{Obstruction Theory}

\begin{theorem}[Obstruction Classes; \ClaimStatusProvedElsewhere]\label{thm:obstructions}
The obstruction to quantizing a $\Pinf$-chiral algebra lies in:
\[
\mathrm{Obs} \in H^3_{\chirPois}(\mathcal{A}_0; \mathcal{A}_0)
\]
the third chiral Poisson cohomology.  The obstruction must vanish at
\emph{each order} in $\hbar$ for quantization to exist to all orders; alternatively,
a formality theorem (as in Theorem~\ref{thm:chiral-kontsevich} below) bypasses the
order-by-order obstruction analysis entirely.
\end{theorem}

\begin{example}[No Obstruction: Affine]\label{ex:no-obs-affine}
For the classical affine Poisson algebra (Example~\ref{ex:classical-affine}):
\[
H^3_{\chirPois}(L\mathfrak{g}^*) = 0
\]
so quantization to $\hat{\mathfrak{g}}_k$ exists for all $k$.
\end{example}

\begin{example}[Obstruction: Anomalous Theories]\label{ex:anomaly}
Certain physical theories have obstructions:
\begin{enumerate}[label=(\roman*)]
\item Chiral WZW model with ``wrong'' level has $\mathrm{Obs} \neq 0$.
\item The obstruction is the \textbf{level quantization condition}: the
      Wess--Zumino class $[H] \in H^3(G, \mathbb{Z})$ requires $k \in \mathbb{Z}$
      for the path integral to be well-defined (arising from $\pi_3(G) \cong \mathbb{Z}$).
\end{enumerate}
\end{example}


\subsection{Anomaly Cancellation}

\begin{theorem}[Green--Schwarz Mechanism; \ClaimStatusProvedElsewhere]\label{thm:green-schwarz}
In string theory, anomalies cancel by introducing a two-form $B$ with modified 
transformation law:
\[
\delta B = \omega_{CS}(\delta A)
\]
where $\omega_{CS}$ is the Chern--Simons form.

At the level of chiral algebras, this modifies the OPE to restore consistency.
\end{theorem}


\subsection{Maurer--Cartan Elements and Deformations}

\begin{definition}[Maurer--Cartan Element]\label{def:mc-element}
For a $\Pinf$-chiral algebra $\mathcal{A}$, a \textbf{Maurer--Cartan element} is 
$\alpha \in \mathcal{A}^1$ (degree 1) satisfying:
\[
d\alpha + \frac{1}{2}\{\alpha, \alpha\} = 0
\]
the Maurer--Cartan equation.
\end{definition}

\begin{theorem}[MC Elements and Quantization; \ClaimStatusProvedElsewhere]\label{thm:mc-quantization}
Quantizations of $\mathcal{A}_0$ correspond to Maurer--Cartan elements in the 
deformation complex:
\[
\mathrm{Def}(\mathcal{A}_0) = (\mathcal{A}_0[[\hbar]] \otimes \mathfrak{g}_{\chirPois}, d + \hbar\{\alpha, -\})
\]
Two quantizations are gauge-equivalent iff their MC elements are related by the 
gauge action.
\end{theorem}

\begin{computation}[MC Equation in Coordinates]\label{comp:mc-coords}
For the classical Heisenberg Poisson algebra with $\{p, q\} = 1$:

The Poisson bivector $\pi = \partial_p \wedge \partial_q$ is the leading-order MC element.  The Kontsevich star product $f \star g = fg + \frac{\hbar}{2}\{f,g\} + O(\hbar^2)$ arises from the formal MC element $\alpha = \hbar\pi + O(\hbar^2)$ in the Hochschild cochain complex.

The MC equation $d\alpha + \frac{1}{2}[\alpha, \alpha] = 0$ at order $\hbar^2$ requires $\frac{1}{2}[\pi, \pi] = 0$, which is the Jacobi identity for the Poisson bracket. Since $\{p,q\} = 1$ is constant, $[\pi,\pi] = 0$ trivially, and $\alpha = \hbar\pi$ is an exact MC element giving the Weyl (Moyal) quantization.
\end{computation}
